\chapter{Présentation}
\label{chap:introduction}

\section{À propos d'AirfoilTools}

\textbf{AirfoilTools} est une bibliothèque Python d'outils dédiés
à l'analyse aérodynamique 2D de profils d'aile. Elle a été extraite
du projet \emph{Axile} (logiciel de conception de parapentes
développé par Nervures) afin de constituer une boîte à outils
réutilisable pour des applications variées : voile, parapente,
hélices, ailes d'avions modèles réduits, etc.

L'application combine :

\begin{itemize}
    \item un \textbf{moteur géométrique} fondé sur les courbes de
        Bézier de degré arbitraire et les splines multi-segment
        (BezierSpline) avec continuités paramétrables (C\textsuperscript{0},
        C\textsuperscript{1}, C\textsuperscript{2}) ;
    \item un \textbf{couplage avec XFoil}, le solveur
        aérodynamique de référence pour les écoulements 2D
        subsoniques visqueux ;
    \item une \textbf{interface graphique interactive} permettant
        l'édition des profils, le pilotage des simulations et la
        visualisation comparative des résultats.
\end{itemize}

\screenshot{01_main_window.png}{1.0}{%
    Vue générale de l'application au démarrage : profils
    NACA~2412 (courant, en bleu) et NACA~0012 (référence,
    en rouge).}

\section{À qui s'adresse ce manuel}

Ce manuel s'adresse aux \textbf{utilisateurs} de l'application
graphique. Il couvre l'ensemble des fonctionnalités accessibles
depuis l'interface, depuis le chargement d'un profil jusqu'à
l'analyse comparative des polaires aérodynamiques.

Les développeurs intéressés par l'API Python de la bibliothèque
trouveront la documentation technique dans le code source
(docstrings au format Sphinx) et dans la documentation générée
(dossier \chemin{docs/}).

\begin{noteinfo}
Aucune connaissance préalable en aérodynamique théorique
n'est requise pour utiliser l'application. En revanche, une
compréhension générale des grandeurs aérodynamiques (coefficient
de portance $C_L$, de traînée $C_D$, finesse $C_L/C_D$,
nombre de Reynolds, angle d'incidence $\alpha$) facilite
l'interprétation des résultats.
\end{noteinfo}

\section{Architecture fonctionnelle}

L'application est structurée autour de \textbf{trois onglets}
correspondant à autant d'étapes logiques :

\begin{enumerate}
    \item \menu{Profils} --- chargement, édition et comparaison des
        profils géométriques. Deux profils sont gérés simultanément :
        un \emph{profil courant} (en bleu) modifiable, et un
        \emph{profil de référence} (en rouge) qui sert de
        comparaison.
    \item \menu{Paramétrage XFoil} --- saisie des paramètres de
        simulation : plage de Reynolds, balayage en incidence,
        paramètres de couche limite, options de panneaux.
    \item \menu{Résultats} --- présentation graphique des résultats
        sous forme de \emph{grille configurable de cellules
        d'analyse}, chacune affichant un tracé indépendant
        (polaire $C_L(\alpha)$, $C_D(\alpha)$, finesse, distribution
        de pression $-C_p(x)$, profil de couche limite, etc.).
\end{enumerate}

\section{Représentation des profils}

AirfoilTools propose deux modes de représentation des profils :

\subsection*{Mode points discrets}

Le profil est défini par une suite de points dans la convention
\textbf{Selig} (parcours du bord de fuite vers l'extrados, puis
le bord d'attaque, puis l'intrados, et retour au bord de fuite).
C'est le format le plus courant des bibliothèques de profils
(par exemple \chemin{airfoiltools.com}).

\subsection*{Mode spline (Bézier multi-segment)}

Le profil est défini par deux \textbf{BezierSpline} (extrados et
intrados), composées chacune de plusieurs segments de Bézier
raccordés en continuité C\textsuperscript{1}. Ce mode permet :

\begin{itemize}
    \item une \emph{édition interactive} par déplacement des
        points de contrôle ;
    \item un \emph{ré-échantillonnage à la demande} avec une
        densité de points librement choisie ;
    \item le \emph{calcul exact de la courbure et des tangentes}
        en tout point ;
    \item le \emph{split d'un segment} en deux à un point arbitraire
        (algorithme de De Casteljau, mathématiquement exact).
\end{itemize}

Le passage d'un mode à l'autre se fait via le menu
\menu{Édition $\rightarrow$ Convertir en Spline} ou par chargement
d'un fichier au format \chemin{.bspl} ou \chemin{.bez}.

\begin{astuce}
Pour bénéficier pleinement des fonctionnalités d'édition
interactive et du tracé de la courbure, il est recommandé
de \textbf{convertir vos profils en mode Spline} dès le
chargement. Les chapitres suivants détaillent cette procédure.
\end{astuce}
